
\section{The geometry of Kullback--Leibler-ball scenarios}
\label{sec:stress}

The main paper keeps the simplest use of this geometry: pricing a published
scenario in nats, which for a pure mean shift is half the squared Mahalanobis
distance. This section gives the full construction of which that use is a
special case, with the proofs that v1 only stated, and a new result on lifting a
portfolio scenario to full dimension.

\subsection{Divergence between two Gaussians}

\begin{proposition}[Closed form]
\label{prop:klgauss}
For $p_1=\mathcal N(\mu_1,\Sigma_1)$ and $p_0=\mathcal N(\mu_0,\Sigma_0)$ on
$\mathbb R^N$, with $\delta=\mu_1-\mu_0$,
\begin{equation}
\label{eq:klgauss}
\KL{p_1}{p_0}=\tfrac12\Big[\underbrace{\delta^\top\Sigma_0^{-1}\delta}_{\text{shift}}
+\underbrace{\operatorname{tr}(\Sigma_0^{-1}\Sigma_1)-N-\log\det(\Sigma_0^{-1}\Sigma_1)}_{\text{deformation}}\Big].
\end{equation}
\end{proposition}

\begin{proof}
$\KL{p_1}{p_0}=\mathbb E_{p_1}[\log p_1-\log p_0]$. Both logarithms are
quadratic; the trace trick of Appendix~\ref{app:gauss} gives
$\mathbb E_{p_1}[(X-\mu_1)^\top\Sigma_1^{-1}(X-\mu_1)]=N$ and
$\mathbb E_{p_1}[(X-\mu_0)^\top\Sigma_0^{-1}(X-\mu_0)]
=\operatorname{tr}(\Sigma_0^{-1}\Sigma_1)+\delta^\top\Sigma_0^{-1}\delta$.
Collect with the $\log\det$ terms.
\end{proof}

The shift term is half the squared Mahalanobis distance: this is the unit in
which the main paper audits published scenarios. The deformation term is
non-negative and zero if and only if $\Sigma_1=\Sigma_0$; writing $\nu_k$ for
the eigenvalues of $\Sigma_0^{-1}\Sigma_1$, it equals
$\tfrac12\sum_k(\nu_k-1-\log\nu_k)$, a sum of non-negative terms, one per
eigendirection.

\subsection{The projected ball and the exchange rate}

\begin{definition}[Entropy-ball scenario]
Given today's conditional law $p_0$ and a severity $\eta>0$, the severity-$\eta$
scenario for a risk functional $\rho$ is
\begin{equation}
\label{eq:ball}
p^\star_\eta=\arg\max_{p\,:\,\KL{p}{p_0}\le\eta}\rho(p).
\end{equation}
\end{definition}

For a linear portfolio $w^\top r$ with mean $m_0$ and standard deviation $s_0$
under $p_0$, set $u=(m-m_0)/s_0$ and $v=s/s_0$. By
Proposition~\ref{prop:klgauss} in one dimension,
\begin{equation}
\label{eq:g}
g(u,v)=\tfrac12\left(u^2+v^2-1-2\log v\right).
\end{equation}

\begin{proposition}[Solution of the projected ball]
\label{prop:ballsol}
For $\rho=-m+c\,s$ with $c\ge0$ (Gaussian VaR and expected shortfall are of this
form), the maximum of $\rho$ subject to $g(u,v)\le\eta$ is attained on the
boundary $g=\eta$ at a point satisfying
\begin{equation*}
u=-t,\qquad v-v^{-1}=c\,t,\qquad t>0 .
\end{equation*}
For $c=0$: $v=1$ and $u=-\sqrt{2\eta}$. \emph{A purely directional scenario of
$\eta$~nats is a move of $\sqrt{2\eta}$~sigmas}; one nat buys
$\sqrt2\approx1.41$~sigma, and $k$~sigmas cost $k^2/2$~nats.
\end{proposition}

\begin{proof}
$g$ is convex on $\mathbb R\times(0,\infty)$ (the sum of $u^2/2$ and
$(v^2-1-2\log v)/2$, whose second derivative is $1+v^{-2}>0$) and the objective
$s_0(-u+cv)$ is linear, so the Karush--Kuhn--Tucker conditions are sufficient.
The objective has no interior maximum, so the constraint binds. The Lagrangian
gives $-1=\lambda u$ and $c=\lambda(v-v^{-1})$ with $\lambda>0$; setting
$t=1/\lambda$ yields the stated system, and $t$ is pinned down by
$g(-t,v(t))=\eta$. For $c=0$, $v-v^{-1}=0$ forces $v=1$ and $g=u^2/2=\eta$.
\end{proof}

\subsection{From the portfolio to full dimension}

The ball~\eqref{eq:ball} is solved on the portfolio's law, but a stress test must
say how each asset moves. Two facts settle the question.

\begin{proposition}[Data processing]
\label{prop:dpi}
For any law $p$ on $\mathbb R^N$, with $p^w$ and $p_0^w$ the laws of $w^\top r$
under $p$ and $p_0$,
$\KL{p^w}{p_0^w}\le\KL{p}{p_0}$.
\end{proposition}

This is the data-processing inequality \citep{cover2006} applied to the map
$r\mapsto w^\top r$. The \emph{ambient} price of a scenario (in full dimension)
is therefore always at least its \emph{portfolio} price (in projection). The
classical-bundle table of v1 bears this out line by line ($9.97$ against
$7.24$~nats for the standard bundle).

\begin{proposition}[Minimal lift]
\label{prop:lift}
Let $p_0=\mathcal N(\mu_0,\Sigma_0)$, $m_0=w^\top\mu_0$, $s_0^2=w^\top\Sigma_0w$,
and a target $(m,s)$ with $s>0$. Among Gaussians $\mathcal N(\mu_1,\Sigma_1)$
with $w^\top\mu_1=m$ and $w^\top\Sigma_1w=s^2$, the divergence $\KL{p_1}{p_0}$
is minimised by
\begin{equation}
\label{eq:lift}
\mu_1=\mu_0+u\,s_0\,\beta,\qquad
\Sigma_1=\Sigma_0+(v^2-1)\,s_0^2\,\beta\beta^\top,
\qquad\beta=\frac{\Sigma_0w}{s_0^2},
\end{equation}
and the minimum equals exactly $g(u,v)$, the portfolio price.
\end{proposition}

\begin{proof}
By Proposition~\ref{prop:dpi}, every candidate costs at least the price of its
projection, which is $g(u,v)$ since the projection of a Gaussian is Gaussian
with mean $m$ and standard deviation $s$. It therefore suffices to check that
\eqref{eq:lift} attains this bound. We have $w^\top\beta=1$, so
$w^\top\mu_1=m_0+us_0=m$ and $w^\top\Sigma_1w=s_0^2+(v^2-1)s_0^2=s^2$. In
\eqref{eq:klgauss}: $\delta^\top\Sigma_0^{-1}\delta=u^2s_0^2\,\beta^\top\Sigma_0^{-1}\beta=u^2$;
$\operatorname{tr}(\Sigma_0^{-1}\Sigma_1)=N+(v^2-1)s_0^2\,\beta^\top\Sigma_0^{-1}\beta=N+v^2-1$;
and by the matrix determinant lemma,
$\det(\Sigma_0^{-1}\Sigma_1)=1+(v^2-1)s_0^2\,\beta^\top\Sigma_0^{-1}\beta=v^2$.
Hence $\KL{p_1}{p_0}=\tfrac12(u^2+v^2-1-2\log v)=g(u,v)$. The matrix $\Sigma_1$ is
positive definite for every $v>0$.
\end{proof}

\paragraph{Asset-by-asset reading.} The vector $\beta$ collects the assets'
betas on the portfolio, $\beta_i=\operatorname{Cov}(r_i,w^\top r)/s_0^2$. Write
$r=\beta f+e$ with $f=w^\top r$ and $e$ uncorrelated with $f$ under $p_0$. The
scenario~\eqref{eq:lift} multiplies the standard deviation of the portfolio
factor $f$ by $v$, shifts its mean by $u$~sigmas, and leaves the law of $e$
unchanged. Each asset's variance rises by $(v^2-1)s_0^2\beta_i^2$ and its mean
moves by $us_0\beta_i$: high-beta assets take the shock, and the correlations
among them rise. It is a one-factor scenario, and it is the least informative of
all scenarios that reach the target severity.

\begin{remark}[Correction to v1]
arXiv~v1 identified this minimum through its first-order conditions,
$\Sigma_1^{-1}=\Sigma_0^{-1}-\lambda ww^\top$, and judged it ``not interpretable
asset by asset''. By the Sherman--Morrison formula, this rank-one perturbation of
the precision matrix \emph{is}~\eqref{eq:lift}, a rank-one perturbation of the
covariance along $\Sigma_0w$, and it reads as above. Since v1 never uses it, no
number depends on it.
\end{remark}

\begin{corollary}[Cost of uniform dilation]
\label{cor:uniform}
Uniform dilation $\Sigma_1=v^2\Sigma_0$, the v1 default, has a deformation
component equal to $N\cdot\tfrac12(v^2-1-2\log v)$, that is $N$~times that of the
minimal lift.
\end{corollary}

\begin{proof}
Every eigenvalue of $\Sigma_0^{-1}\Sigma_1$ equals $v^2$; apply the spectral form
of the deformation term.
\end{proof}

For the ten-year scenario on the main panel ($v=1.63$, $N=20$), the variance
component is $0.34$~nat at the minimum and $20\times0.34\approx6.8$ under uniform
dilation, which is the figure in the v1 composition table. The corollary
explains the v1 finding that the compositions calibrated on H1 ($2.88$~nats) and
refined by H8 ($0.71$) cost \emph{less} than the default: the default deforms
every eigendirection equally, whereas a composition that concentrates the shock
on the common mode moves towards the one-factor scenario~\eqref{eq:lift}, whose
total cost is $0.41$~nat including the mean shift. That uniform dilation costs
$N$~times the minimum in variance is a general result; that the calibrated
compositions fall between the two remains a finding on these panels, but it
becomes expected rather than surprising.

\subsection{Composing a given severity}

\begin{proposition}[One severity, many compositions]
\label{prop:composition}
If $w^\top\Sigma_aw=w^\top\Sigma_bw=s^2$, then for every $\alpha\in[0,1]$,
$\Sigma_\alpha=(1-\alpha)\Sigma_a+\alpha\Sigma_b$ satisfies
$w^\top\Sigma_\alpha w=s^2$.
\end{proposition}

\begin{proof}
$w^\top\Sigma w$ is linear in $\Sigma$.
\end{proof}

Portfolio severity therefore does not constrain the composition. The ambient
price, for its part, is convex in $\Sigma_1$ (the $-\log\det$ term is convex and
the rest is linear), so along a mixture it is bounded by the mixture of prices,
which guarantees that an intermediate composition never costs more than the
dearer of its two endpoints.

\paragraph{A physical ceiling on correlation.} With marginal volatilities fixed,
$w^\top\Sigma w$ is maximised over valid correlation matrices at
$\rho_{ij}=1$, where it equals $(\sum_iw_i\sigma_i)^2$. On the main panel this
ceiling is only about $1.5$~times today's volatility: beyond the five-year rung,
no rise in correlation alone suffices, and some rise in volatilities is
unavoidable.

\subsection{Calibrating the radius, and the limit of that calibration}

The radius is read off the information flow of the projection,
$I^{(h)}_t=\KL{p^w_t}{p^w_{t-h}}$, whose quantiles (on non-overlapping 21-day
blocks) convert a return period into a radius. Reading the radius off the
$N$-dimensional flow overstates it by an order of magnitude, by
Proposition~\ref{prop:dpi}. Calibrated this way, the v1 ladder gives a ten-year
volatility multiplier of $1.63$, while the ten-year quantile of the realised
one-month volatility rise is $3.2$. The ladder is a coherent, reproducible rule;
its return periods are not validated. This is one of the reasons the main paper
keeps the audit (pricing a given scenario) and drops generation (building a
scenario with a given return period).

\subsection{Reverse stress testing}

Asking for the cheapest scenario that produces a given loss $L$ inverts the ball.
In the projected Gaussian setting, Proposition~\ref{prop:ballsol} reads
backwards. In the non-parametric setting, the solution is the exponential
tilting of the sample \citep{csiszar1975}, whose effective sample size collapses
as the radius grows: at $9$~nats it is $1.0$ on the main panel, meaning that the
tilted law is nothing but the worst observed day. Any implementation should
report this effective size and switch to the Gaussian ball as soon as it drops
below a few dozen.
